\section{What the detector families have reported}\label{sec:literature}

\Cref{tab:families} places the families on the two numbers of \cref{sec:scaling}. The first column group is the geometry, the
last three are what was reported and how it was handled. All lever arms are derived from the published geometry and
diffusion coefficients, since no source states the quantity; where a diffusion coefficient had to be assumed it is marked.

\begin{table}[t]
  \centering\footnotesize
  \caption{Detector families against the two numbers that control the single-time bias: the strips in a full-drift column at
  $\tan\theta = 0.3$ and the transverse diffusion at full drift in units of the pitch. $\Lambda_{0.3} = \sigma_T(L)/(L\tan\theta)$
  is the lever arm. Derived values rest on the cited geometry; the LMU diffusion coefficients are assumed equal to the NSW value.}
  \label{tab:families}
  \begin{tabular}{@{}p{3.1cm}cccc>{\raggedright\arraybackslash}p{3.0cm}>{\raggedright\arraybackslash}p{2.9cm}@{}}
    \toprule
    Family (examples) & $L$ [mm] & $p$ [mm] & $N_{0.3}$ & $\sigma_T(L)/p$ & Time estimator; reported bias & Treatment \\
    \midrule
    MX17 chamber A (this work) & 28 & 0.78 & 10.8 & 0.75 & leading edge, peak, template, centroid: ratio 1.19 (X), 1.37 (Y) at the sample median & forward fit of the waveforms \\
    ATLAS NSW Micromegas~\cite{Iodice2015,Ntekas2016,Scholer2021} & 5 & 0.425--0.45 & 3.4 & 0.6 & APV25 rise time, VMM peak; \ang{20}$\to$\ang{23}; sharing-off simulation still $+\ang{7}$ & edge filters, position polynomial, angle cuts \\
    LMU floating strip~\cite{Bortfeldt2014,Bortfeldt2015} & 6 & 0.5 & 3.6 & 0.56 & Fermi fit; \ang{10}$\to$\ang{15} measured, \ang{19} in LTspice & angle-dependent rescale \\
    LMU GEM~\cite{Flierl2018} & 8 & 0.4 & 6 & 0.82 & Fermi fit; \ang{15}$\to$\ang{19}, \ang{30}$\to$\ang{32} & 15--20\,\% neighbour subtraction \\
    BESIII CGEM~\cite{Alexeev2019,Farinelli2019} & 5 & 0.65 & 2.3 & 0.38 & APV25; ``tilt towards larger angles'' & edge weighting, \SI{180}{\micro\metre} shift \\
    T2K ERAM~\cite{T2K2022,T2K2023} & 970 & 10--11 & 27 & 0.3 & time of maximum; RC delay, no angle bias reported & RC and gain maps; no time ladder \\
    ALICE, STAR pad TPCs~\cite{Alme2010,Anderson2003} & 2100--2500 & 5--7 & 110 & 0.6--0.9 & charge-weighted time per row; resolution terms only & drift scale from lasers and silicon \\
    LArTPCs~\cite{Adams2018a,Adams2018b,ProtoDUNE2020,ICARUS2024} & 1500--3600 & 3--4.8 & 230 & 0.3--0.8 & Gaussian peak after 2-D deconvolution; width and charge vs angle, no time bias found & deconvolution in wire and time \\
    Directional and short-track TPCs~\cite{MIMAC2020,BEAST2019,MAIKo2018} & 10--250 & 0.05--2 & few & 0.5--8 & threshold, TOT midpoint, slice barycentre; ``canoe'' shapes, slice ends deviate & none; resolution \ang{1}--\ang{60} \\
    GridPix~\cite{GridPix2018} & 20 & 0.055 & 110 & 8 & time of each electron & nothing to aggregate \\
    \bottomrule
  \end{tabular}
\end{table}

\subsection{Strip micro-TPCs: the bias is a family trait}

Every strip $\mu$TPC that published a comparison with a set angle reports the track too steep, with the excess shrinking as
the angle grows. Bortfeldt measured $+\ang{5}$ at \ang{10} on a floating-strip Micromegas and predicted $+\ang{9}$ with a
parameter-free LTspice model of the strip capacitances, which also reproduced the fall with angle (\ang{4.3}, \ang{2.0}, \ang{0.6}
at \ang{20}, \ang{30}, \ang{40})~\cite{Bortfeldt2014}; the journal version names two effects, ``signals couple capacitively onto
neighboring strips'' so that ``measured signal timing is shifted towards earlier values for the strips hit by electrons created
close to the cathode'', and ``for strips at the edges of the strip cluster, the ionization electrons position is not calculated
correctly''~\cite{Bortfeldt2015}. ATLAS test chambers at \ang{20} peak near \ang{23}~\cite{Iodice2015}, ``understood to be due to
the capacitive induction of the signal on neighboring strips'', a formulation the BESIII GEM papers repeat nearly verbatim while
adding that ``the transverse diffusion plays another important role''~\cite{Alexeev2019}. Flierl's GEM read \ang{19} (most probable)
at \ang{15} and \ang{32} (mean) at \ang{30}~\cite{Flierl2018}. In the units of \cref{eq:scaling} these are tan ratios of 1.52, 1.17, 1.28
and 1.08 at $N$ of 2.1, 4.5, 5.4 and 11.5 (\cref{fig:universal}b), and $(\text{ratio}-1)N^2$ lies between 2 and 11 for all of
them.

The attribution rests on two roots, Bortfeldt's circuit model~\cite{Bortfeldt2014} and the NTUA/ATLAS Maxwell--Simplorer model
with an edge-strip charge correlation~\cite{Ntekas2016}; later papers cite or restate them~\cite{Sessa2014,Vanadia2015,Rinnagel2019,Alexeev2019,Farinelli2019}.
Neither root excludes the alternatives. Bortfeldt's model has coupling and the edge-position effect but no diffusion, and
gives no decomposition; Ntekas's thesis states that the gas-only Monte Carlo ``without the fully equivalent electrical circuit''
still shows a bias, ``purely geometrical coming from the incomplete overlap of the outer cluster strips with the particle's
trace inside the MM gas gap combined with the diffusion of the ionisation electrons'', of \ang{0.2} at \ang{30} rising to \ang{4} at
\ang{10}~\cite{Ntekas2016}. Scholer's NSW digitisation, with diffusion on and charge sharing switched off, peaks \ang{7} high for
tracks above \ang{25}, an offset he attributes to ``the positions of the first and last ionizations above the corresponding
strips'' and leaves to ``further investigation''; with sharing on the peak moves higher, and his own test beam at \ang{28} shows no
shift at all~\cite{Scholer2021}. That test beam is the one published measurement that lies below the no-coupling
toy of \cref{app:toy}, which on the NSW-like geometry with the VMM peak time gives a tangent ratio of 1.10 at \ang{28}, a
peak about \ang{2} high; his own simulation, \ang{7} high, misses it by more. Two differences are not checked: the quoted value is
the peak of a distribution with a long tail towards larger angles (resolution \ang{7.1} about the peak), not the median our
ratios use, and a $6\times$ r.m.s.\ threshold that drops the faint end strips acts as an edge cut (\cref{sec:remedies}).
Klitzner, with floating strips in the fast gas Ne:CF$_4$, finds the angle too large at drift fields below 300~V/cm, but
on one layer only: ``all other layers reconstruct a slightly smaller angle than expected''~\cite{Klitzner2016}. These are
the two published cases where the bias is absent or reversed; neither author explains them, and we cannot either. Herrmann, on a full-size NSW module, adopts a 30\,\% coupling correction from Lösel's
simulation~\cite{Loesel2017} and Flierl's analysis, and states the mechanism in strip-count terms: without the correction ``no
mean-time difference is in the range expected from the drift velocity'', and ``the time difference is biased for lower
inclination due to the low number of affected strips''~\cite{Herrmann2019}. The only other mechanism statement we found is a
2024 conference poster by Fang et al., who derive a per-strip time correction from transverse and longitudinal diffusion on the
argument that the first of several primary electrons sets the strip time~\cite{Fang2024}; it concerns the segment midpoint and
has no journal version.

Two facts stand out. No source varied one parameter and tracked the bias: pitch, strip length, drift field and electronics are
confounded in every comparison, and a drift-field scan changes the diffusion, the velocity, the time per strip and the signal
to noise together. And no source had a per-track external angle; the set angle of a chamber is known to about a degree, which at
\ang{10} is 10\,\% in the tangent, the size of the effect under study.

\subsection{Pad TPCs and liquid-argon TPCs: the same spread, a longer column, another estimator}

The large pad TPCs have a transverse diffusion at full drift comparable to their pad width (ALICE 0.6--0.9, STAR 0.5--1.0 in
units of the pad~\cite{Alme2010,Anderson2003}), the same as MX17's 0.75. What differs is the column: 45 to 150 rows against 3
to 13 strips, so the end strips are a fraction of a per cent of the points. They also read a charge-weighted time, STAR
``weighting the average by the amount of charge collected in each bucket''~\cite{Anderson2003} and ALICE a centre of gravity over
time bins~\cite{Alme2010}, which has no interior compression. Their angle-dependent terms (the angular pad and angular wire effects,
$\tan^2\alpha$ in the resolution, threshold losses) are variances, not slope biases, and their drift scale is fixed by lasers
and silicon to 0.1\,\%. A multiplicative compression of the drift coordinate would be absorbed into the effective drift
velocity by a track-based calibration~\cite{Alme2010}; at the 0.1\,\% our scaling predicts for a 1.3~m column there is nothing to
absorb.

The liquid-argon TPCs are the closest relatives of the forward model. MicroBooNE writes the wire signal as the charge
convolved with a field and electronics response in wire and time and solves it as a two-dimensional deconvolution, with
Wiener-type filters, because the long-range induction puts a signal on wires that collect nothing~\cite{Adams2018a,Adams2018b};
ProtoDUNE-SP uses the same form~\cite{ProtoDUNE2020}. Their kernel is the induced current, bipolar and calculated; the transverse
diffusion stays in the deconvolved image. MicroBooNE states the mechanism for the pulse width, adjacent clouds that ``begin to
overlap as they spread in the $yz$-plane under the influence of $D_T$, causing additional $\sigma_t$ smearing'' that grows with the
angle to the wire plane~\cite{MicroBooNE2021}, and ICARUS tunes its wire response in angle bins~\cite{ICARUS2024}. A bias of the
hit time or of the local slope has not been reported; with columns of hundreds of wires it would be below a per mille, and a
3~cm segment at full drift, where the lever arm is as short as MX17's, has not been tested against an external slope.

\subsection{Directional and short-track detectors: the column but no reference}

The keV-to-MeV recoil detectors have short tracks on fine channels and one time per channel (NEWAGE, BEAST, D3, MAIKo,
MIGDAL's strip plane), or a barycentre per time slice (MIMAC, ACTAR)~\cite{MIMAC2020,BEAST2019,MAIKo2018,MIGDAL2023}. Two of them
describe the artefact without quantifying it: BEAST's pixel tracks take a ``canoe shape'' because ``the pixel integration time
is longer than the typical charge cloud size''~\cite{BEAST2019}, and MIMAC's ``initial and final time slices of the track
usually have a larger than average deviation from the track direction''~\cite{MIMAC2020}. Their angle references are
kinematics at \ang{10} to \ang{16} resolution, source pointing with no absolute drift coordinate, or the self-consistency of the two
halves of one track, which a common scale error passes. A few per cent of slope is invisible to all of them. The single-electron
pixel detectors (GridPix) have nothing to aggregate and report only a time walk~\cite{GridPix2018}; their drift scale is
calibrated with a telescope, which would absorb a compression~\cite{GasPix2019}.

\subsection{Drift tubes: the one time per wire is already a functional of the inclination}

The oldest time-projection readout shows the general statement in another guise. The drift time of a wire is the arrival of
the first electrons from the track segment seen by that cell, and its relation to the distance of closest approach depends on
the track's inclination to the cell: the CMS drift tubes calibrate a maximum drift time that ``depends on the track inclination
and on the pattern of cells crossed'' and an effective drift velocity that depends on ``the inclination of the
track''~\cite{CMSDT2010}; the BESIII drift chamber determines separate space-time relations ``for 10 degree steps'' of the entrance
angle because the differences ``grow with the increasing drift times up to as much as several hundred microns''~\cite{BESIIIMDC2014};
Belle II parametrises its relations per layer, side and angle category~\cite{BelleIICDC2019}. These detectors never see the
problem as a bias because the angle-dependent calibration is part of the method. The $\mu$TPC has no such calibration, because
its purpose is to measure the inclination that the calibration would have to know.

\subsection{What survives the survey}

A transparent toy with drift diffusion, the strip geometry and a small induction footprint, and no coupling of any kind,
reproduces the published biases of the NSW-type, floating-strip and GEM chambers to within a factor 1.6 in the excess
$\tan\theta_{\rm rec}/\tan\theta - 1$ on the four measured points (toy over published 1.6, 1.3, 0.8 and 0.6, \cref{tab:toydet}), with no
free parameter beyond an assumed footprint and diffusion; the circuit model's predictions at \ang{20} to \ang{40} and the diffusion-only toy differ by
less than 0.05 in the tangent ratio. The angle, in other words, never contained the information that would have identified
the coupling.

The bias is not new, the model equation is not new, and the per-channel waveform fit has relatives in T2K's RC maps and in the
LArTPC deconvolution. What had not been done is to measure the compression per track against an independent reference, to
separate the geometric end term from the lateral spreads and the spreads from each other, and to fit the drift-depth charge
profile of each track inside the convolution model. Those are the contributions of \cref{sec:measurements,sec:remedies}, and the
companion paper~\cite{paperII}.
